% GENERATED FILE. Edit canonical lessons, then run npm run generate:books.
% canonical-source-hash: fnv1a64:05698ab97f891afa
% canonical-lessons: SA-C398-ghattah, SA-C398-avagahate, SA-C398-sangamah, SA-C398-pavanah, SA-C398-proksati

\chapter{At the River Steps}
\label{ch:sa-at-the-river-steps}
\hlchaptermodality{398}

\begin{chapteropening}
\textbf{By the end of this chapter:} \emph{I can say where pilgrims bathe and what they do there.}

Complete the last lesson of chapter 398: I can say where pilgrims bathe and what they do there.
\end{chapteropening}

\section[\texorpdfstring{\sk{घट्टः} (ghaṭṭaḥ)}{घट्टः (ghaṭṭaḥ)}]{\sk{घट्टः} (ghaṭṭaḥ) — a bathing ghat; river steps}

\label{lesson:SA-C398-ghattah}



\begin{quote}
Before the new one: say the Sanskrit for a hymn of praise, then the Sanskrit for sings the praises of.
\end{quote}

\subsection*{You'll want to know: \sk{घट्टः}}
\textbf{\sk{घट्टः}} — \emph{ghaṭṭaḥ} — \textquotedblleft{}a bathing ghat; river steps\textquotedblright{}.

\begin{quote}
\sk{प्रातः} \sk{घट्टे} \sk{यात्रिकाः} \sk{स्नान्ति।}
\end{quote}

\begin{quote}
\emph{prātaḥ ghaṭṭe yātrikāḥ snānti.}
\end{quote}

\begin{quote}
In the morning pilgrims bathe at the bathing ghat.
\end{quote}

A masculine noun.

\begin{cousinweb}
\textbf{\sk{घट्टः}} (\emph{ghaṭṭaḥ}) is the word MW gives for a landing-place, a bathing place or steps by a river, with no derivation; CDIAL notes that a Dravidian origin has been suggested (T. Burrow), so its origin is uncertain; it lives on as Hindi \emph{ghāṭ}, \textquotedblleft{}a bathing or landing place\textquotedblright{} (MW; CDIAL 4414).
\end{cousinweb}

\begin{grammarlens}[title={What you've built}]
One more word for getting around.
\end{grammarlens}

\subsection*{Your turn}
\begin{itemize}
\raggedright
  \item \emph{Say it:} \emph{ghaṭṭaḥ}
  \item \emph{Say it:} \emph{ghaṭṭaḥ}, once more
  \item \emph{Say it:} \emph{prātaḥ ghaṭṭe yātrikāḥ snānti.} — In the morning pilgrims bathe at the bathing ghat.
  \item \emph{Recall:} say the Sanskrit for a hymn of praise, then the Sanskrit for sings the praises of, then say \emph{ghaṭṭaḥ} again
\end{itemize}

\begin{tcolorbox}[breakable,colback=teal!4,colframe=teal!35!black,title={Before you move on}]
What does \sk{घट्टः} mean? (\textquotedblleft{}A bathing ghat; river steps\textquotedblright{}.) Say it once more.
\end{tcolorbox}
\section[\texorpdfstring{\sk{अवगाहते} (avagāhate)}{अवगाहते (avagāhate)}]{\sk{अवगाहते} (avagāhate) — plunges into; bathes in}

\label{lesson:SA-C398-avagahate}



\begin{quote}
Before the new one: say the Sanskrit for sings the praises of, then the Sanskrit for a bathing ghat.
\end{quote}

\subsection*{You'll want to know: \sk{अवगाहते}}
\textbf{\sk{अवगाहते}} — \emph{avagāhate} — \textquotedblleft{}plunges into; bathes in\textquotedblright{}.

\begin{quote}
\sk{यात्रिकः} \sk{नद्याम्} \sk{अवगाहते।}
\end{quote}

\begin{quote}
\emph{yātrikaḥ nadyām avagāhate.}
\end{quote}

\begin{quote}
The pilgrim plunges into the river.
\end{quote}

Past (he …): \sk{सः} \sk{नद्याम्} \sk{अवगाढवान्}, \emph{saḥ nadyām avagāḍhavān}.

\begin{cousinweb}
\textbf{\sk{अवगाहते}} (\emph{avagāhate}) is built from \emph{ava}, \textquotedblleft{}down\textquotedblright{}, and the root √gāh, \textquotedblleft{}to plunge, to dive into\textquotedblright{}; MW gives \emph{ava-gāh} as to plunge into or bathe in, with the participle \emph{avagāḍha} (MW).
\end{cousinweb}

\begin{grammarlens}[title={What you've built}]
One more word for getting around.
\end{grammarlens}

\subsection*{Your turn}
\begin{itemize}
\raggedright
  \item \emph{Say it:} \emph{avagāhate}
  \item \emph{Say it:} \emph{avagāhate}, once more
  \item \emph{Say it:} \emph{yātrikaḥ nadyām avagāhate.} — The pilgrim plunges into the river.
  \item \emph{Recall:} say the Sanskrit for sings the praises of, then the Sanskrit for a bathing ghat, then say \emph{avagāhate} again
\end{itemize}

\begin{tcolorbox}[breakable,colback=teal!4,colframe=teal!35!black,title={Before you move on}]
What does \sk{अवगाहते} mean? (\textquotedblleft{}Plunges into; bathes in\textquotedblright{}.) Say it once more.
\end{tcolorbox}
\section[\texorpdfstring{\sk{संगमः} (saṅgamaḥ)}{संगमः (saṅgamaḥ)}]{\sk{संगमः} (saṅgamaḥ) — a confluence of rivers}

\label{lesson:SA-C398-sangamah}



\begin{quote}
Before the new one: say the Sanskrit for a bathing ghat, then the Sanskrit for plunges into.
\end{quote}

\subsection*{You'll want to know: \sk{संगमः}}
\textbf{\sk{संगमः}} — \emph{saṅgamaḥ} — \textquotedblleft{}a confluence of rivers\textquotedblright{}.

\begin{quote}
\sk{संगमे} \sk{जलं} \sk{शीतलम्} \sk{अस्ति।}
\end{quote}

\begin{quote}
\emph{saṅgame jalaṃ śītalam asti.}
\end{quote}

\begin{quote}
The water at the confluence is cool.
\end{quote}

A masculine noun.

\begin{cousinweb}
\textbf{\sk{संगमः}} (\emph{saṅgamaḥ}) is built from \emph{sam}, \textquotedblleft{}together\textquotedblright{}, and the root √gam, \textquotedblleft{}to go\textquotedblright{}: a going together; MW gives \emph{saṃgama} as a meeting and as the confluence of rivers (MW).
\end{cousinweb}

\begin{grammarlens}[title={What you've built}]
One more word for getting around.
\end{grammarlens}

\subsection*{Your turn}
\begin{itemize}
\raggedright
  \item \emph{Say it:} \emph{saṅgamaḥ}
  \item \emph{Say it:} \emph{saṅgamaḥ}, once more
  \item \emph{Say it:} \emph{saṅgame jalaṃ śītalam asti.} — The water at the confluence is cool.
  \item \emph{Recall:} say the Sanskrit for a bathing ghat, then the Sanskrit for plunges into, then say \emph{saṅgamaḥ} again
\end{itemize}

\begin{tcolorbox}[breakable,colback=teal!4,colframe=teal!35!black,title={Before you move on}]
What does \sk{संगमः} mean? (\textquotedblleft{}A confluence of rivers\textquotedblright{}.) Say it once more.
\end{tcolorbox}
\section[\texorpdfstring{\sk{पावनः} (pāvanaḥ)}{पावनः (pāvanaḥ)}]{\sk{पावनः} (pāvanaḥ) — purifying; holy}

\label{lesson:SA-C398-pavanah}



\begin{quote}
Before the new one: say the Sanskrit for plunges into, then the Sanskrit for a confluence.
\end{quote}

\subsection*{You'll want to know: \sk{पावनः}}
\textbf{\sk{पावनः}} — \emph{pāvanaḥ} — \textquotedblleft{}purifying; holy\textquotedblright{}.

\begin{quote}
\sk{नद्याः} \sk{जलं} \sk{पावनम्} \sk{अस्ति।}
\end{quote}

\begin{quote}
\emph{nadyāḥ jalaṃ pāvanam asti.}
\end{quote}

\begin{quote}
The river's water is purifying.
\end{quote}

\begin{cousinweb}
\textbf{\sk{पावनः}} (\emph{pāvanaḥ}) comes from the root √pū, \textquotedblleft{}to make clean, to purify\textquotedblright{}; MW gives \emph{pāvana} as purifying, pure, holy (MW).
\end{cousinweb}

\begin{grammarlens}[title={What you've built}]
One more word for getting around.
\end{grammarlens}

\subsection*{Your turn}
\begin{itemize}
\raggedright
  \item \emph{Say it:} \emph{pāvanaḥ}
  \item \emph{Say it:} \emph{pāvanaḥ}, once more
  \item \emph{Say it:} \emph{nadyāḥ jalaṃ pāvanam asti.} — The river's water is purifying.
  \item \emph{Recall:} say the Sanskrit for plunges into, then the Sanskrit for a confluence, then say \emph{pāvanaḥ} again
\end{itemize}

\begin{tcolorbox}[breakable,colback=teal!4,colframe=teal!35!black,title={Before you move on}]
What does \sk{पावनः} mean? (\textquotedblleft{}Purifying; holy\textquotedblright{}.) Say it once more.
\end{tcolorbox}
\section[\texorpdfstring{\sk{प्रोक्षति} (prokṣati)}{प्रोक्षति (prokṣati)}]{\sk{प्रोक्षति} (prokṣati) — sprinkles over; sprinkles water on}

\label{lesson:SA-C398-proksati}



\begin{quote}
Before the new one: say the Sanskrit for a confluence, then the Sanskrit for purifying.
\end{quote}

\subsection*{You'll want to know: \sk{प्रोक्षति}}
\textbf{\sk{प्रोक्षति}} — \emph{prokṣati} — \textquotedblleft{}sprinkles over; sprinkles water on\textquotedblright{}.

\begin{quote}
\sk{पुत्रः} \sk{पुष्पेषु} \sk{जलं} \sk{प्रोक्षति।}
\end{quote}

\begin{quote}
\emph{putraḥ puṣpeṣu jalaṃ prokṣati.}
\end{quote}

\begin{quote}
The son sprinkles water on the flowers.
\end{quote}

Past (he …): \sk{सः} \sk{पुष्पेषु} \sk{जलं} \sk{प्रोक्षितवान्}, \emph{saḥ puṣpeṣu jalaṃ prokṣitavān}.

\begin{cousinweb}
\textbf{\sk{प्रोक्षति}} (\emph{prokṣati}) is built from \emph{pra}, \textquotedblleft{}forth\textquotedblright{}, and the root √ukṣ, \textquotedblleft{}to sprinkle\textquotedblright{}; MW gives \emph{prokṣ} as to sprinkle upon or to purify by sprinkling, with the participle \emph{prokṣita} (MW).
\end{cousinweb}

\begin{grammarlens}[title={What you've built}]
One more word for getting around.
\end{grammarlens}

\subsection*{Your turn}
\begin{itemize}
\raggedright
  \item \emph{Say it:} \emph{prokṣati}
  \item \emph{Say it:} \emph{prokṣati}, once more
  \item \emph{Say it:} \emph{putraḥ puṣpeṣu jalaṃ prokṣati.} — The son sprinkles water on the flowers.
  \item \emph{Recall:} say the Sanskrit for a confluence, then the Sanskrit for purifying, then say \emph{prokṣati} again
\end{itemize}

\begin{tcolorbox}[breakable,colback=teal!4,colframe=teal!35!black,title={Before you move on}]
What does \sk{प्रोक्षति} mean? (\textquotedblleft{}Sprinkles over; sprinkles water on\textquotedblright{}.) Say it once more.
\end{tcolorbox}
